\newpage
\subsection{Schulungen}

NIS2 fordert regelmäßige IT-Sicherheitsschulungen für alle Mitarbeitenden; die Geschäftsleitung muss zusätzlich an Schulungen zum IT-Risikomanagement teilnehmen. \textit{Regelmäßig} im Sinne von NIS2 heißt mindestens alle drei Jahre. Empfehlenswert sind jedoch jährliche Schulungen, da sich Bedrohungen und Angriffsmethoden ständig weiterentwickeln.

Ein Schulungskonzept sollte klären: Wer schult? Welche Themen sind relevant? Welche Formate werden eingesetzt? Die Inhalte können je nach Organisation und Zielgruppe variieren; manche Mitarbeitende benötigen tätigkeitsbedingt vertiefende Schulungen. Ob intern geschult oder an einen externen Dienstleister ausgelagert wird, gibt NIS2 nicht vor.

\subsubsection{Schulungsformate}

Je nach Unternehmensgröße und Ressourcen eignen sich unterschiedliche Formate:

\begin{itemize}
    \item \textbf{Präsenzschulungen}: Klassische Schulungen vor Ort, die Interaktion und Rückfragen ermöglichen
    \item \textbf{E-Learning}: Flexible Online-Kurse, die Mitarbeitende zeitunabhängig absolvieren können
    \item \textbf{Micro-Learning}: Kurze Lerneinheiten (5–10 Minuten), die regelmäßig bereitgestellt werden
    \item \textbf{Phishing-Simulationen}: Kontrollierte Test-Phishing-Mails, um das Erkennen von Angriffen zu trainieren
    \item \textbf{Awareness-Kampagnen}: Plakate, Newsletter oder Intranet-Beiträge zu aktuellen Bedrohungen
\end{itemize}

Am wirksamsten ist oft eine Kombination (Blended Learning), da die Themen über verschiedene Medien vermittelt werden und so präsent bleiben.
\newpage
\subsubsection{Schulungsinhalte}

\textbf{Eigenes Sicherheitskonzept}
\begin{itemize}
    \item Struktur und Abläufe des unternehmensweiten IT-Sicherheitskonzepts
    \item Rollen und Verantwortlichkeiten (z.\,B. IT-Sicherheitsbeauftragter, Datenschutzkoordinator)
    \item Meldewege für Sicherheitsvorfälle und Datenschutzverstöße
    \item Umgang mit privaten Geräten (BYOD) im Unternehmensnetzwerk
\end{itemize}

\textbf{Cyberhygiene}
\begin{itemize}
    \item Social Engineering: Telefonbetrug, Pretexting, Baiting, Tailgating
    \item Malware: Einfallstore, Typen (Ransomware, Trojaner, Spyware), Auswirkungen
    \item Phishing: Massen-Phishing, Spear-Phishing, CEO-Fraud, Smishing, Vishing
    \item Passwortsicherheit und Zwei-Faktor-Authentifizierung (2FA/MFA)
    \item Umgang mit Wechselmedien: USB-Sticks, externe Festplatten
    \item Verschlüsselung und sichere Verbindungen: HTTPS, VPN, E-Mail-Verschlüsselung
    \item Verhalten in öffentlichen WLANs und unsicheren Netzwerken
    \item Absicherung mobiler Endgeräte: Smartphones, Tablets, Notebooks
    \item Sauberer Arbeitsplatz (Clean Desk Policy) und gesperrte Bildschirme
    \item Umgang mit Cloud-Diensten (Filesharing, Speicherung sensibler Daten)
    \item Erkennen und Melden von verdächtigen Aktivitäten
    \item Aktuelle Trends und neueste Bedrohungen
\end{itemize}

\textbf{Datenschutz}
\begin{itemize}
    \item Grundprinzipien: Zweckbindung, Datenminimierung, Vertraulichkeit
    \item Kommunikation mit Kunden, Dienstleistern und externen Partnern
    \item Vermeidung unbedachter Veröffentlichungen (z.\,B. Social Media)
    \item Umgang mit personenbezogenen und sensiblen Daten
    \item Datenschutz bei mobiler Arbeit und Homeoffice
    \item Sichere Aufbewahrung und Entsorgung sensibler Informationen
\end{itemize}
\newpage

Diese Inhalte dienen als Orientierung; je nach Organisation können weitere Themen relevant sein. Hilfreich ist es, im Team relevante Themen zu besprechen, vergangene Sicherheitsvorfälle zu analysieren und aktuelle Bedrohungen zu recherchieren. Das Schulungskonzept sollte regelmäßig überprüft und angepasst werden.

NIS2 fordert die Teilnahme aller Mitarbeitenden. In größeren Unternehmen ist eine vollständige Abdeckung nicht immer realistisch, eine möglichst hohe Teilnahmequote sollte dennoch angestrebt werden.

Folgende Punkte müssen dokumentiert werden:
\begin{itemize}
    \item Datum und Dauer der Schulung
    \item Schulungsinhalte
    \item Teilnehmende (mit Unterschrift oder elektronischem Nachweis)
    \item Schulungsformat (Präsenz, E-Learning etc.)
    \item Bei der Geschäftsleitung: gesonderte Dokumentation der Risikomanagement-Schulung
\end{itemize}

Neue Mitarbeitende sollten idealerweise in den ersten Wochen eine IT-Sicherheitsschulung erhalten, die mindestens die Grundlagen vermittelt: Meldewege bei Vorfällen, Passwortrichtlinien, Umgang mit E-Mails und Phishing sowie die grundlegenden Sicherheitsregeln. Die vollständige Awareness-Schulung kann später folgen, aber nicht später als nach drei Monaten.

\subsubsection{Anlassbezogene Schulungen}

Zusätzlich zu den regelmäßigen Schulungen sollte anlassbezogen geschult werden:

\begin{itemize}
    \item Nach Sicherheitsvorfällen im eigenen Unternehmen
    \item Bei neuen oder veränderten Bedrohungslagen (z.\,B. neue Phishing-Kampagnen)
    \item Bei Einführung neuer Systeme oder Prozesse mit Sicherheitsrelevanz
    \item Bei Änderungen der Sicherheitsrichtlinien
\end{itemize}


\newpage
\vspace{1em}
\textbf{Checkliste: Schulungen}
\begin{checklist}
    \item Identifizieren welche Daten in der Organisation verschlüsselt werden müssen
    \item Bei internen Schulungen: Verantwortliche Personen festlegen
    \item Bei externen Schulungen: Anbieter prüfen und auswählen
    \item Bei externen Schulungen: Dienstleister im Rahmen des Lieferantenmanagements bewerten
    \item Schulungskonzept erstellen (Zielgruppen, Inhalte, Formate, Häufigkeit)
    \item Onboarding-Prozess für neue Mitarbeitende festlegen
    \item Methoden zur Erfolgskontrolle definieren (Tests, Phishing-Simulationen)
    \item Prozess zur Dokumentation der Schulungsinhalte, Teilnehmer und Datum der Schulung
    \item Schulungskonzept regelmäßig überprüfen und aktualisieren
\end{checklist}
